\newpage
\pagestyle{empty}

\begin{center}
  \bf\Large 研究生学位论文/实践成果报告使用人工智能工具声明
\end{center}

\vskip 1cm

\normalsize \indent
就本学位论文/实践成果报告的研究与撰写过程，是否使用生成式人工智能辅助工具情况，声明如下（根据实际情况，在以下两项中选择一项打勾）：

\vskip 0.5cm

\noindent【\quad】1.本学位论文/实践成果报告的研究与撰写过程未使用任何生成式人工智能工具（指能够根据提示自动生成文本、代码、图表等内容的大模型及其衍生工具）。

\noindent 所有内容均基于个人研究、文献阅读及实验数据完成，本人承诺符合学术诚信要求，论文/成果中不存在任何人工智能辅助生成内容。违反本声明所产生的责任与后果，由本人独立承担。

\vskip 0.5cm

\noindent【\quad】2.本学位论文/实践成果报告的研究与撰写过程使用了生成式人工智能辅助工具，现按照要求披露如下：

\vskip 0.3cm

\noindent（1）使用的生成式人工智能工具名称及目的

\noindent 工具名称：\underline{\hspace{6cm}}（例如：DeepSeek、文心一言、豆包、智谱、千问等 AI 大模型，以及调用 AI 大模型的智能体工具等）

\noindent 具体用途：\underline{\hspace{6cm}}（例如：文献检索、语言润色、数据处理等）

\vskip 0.3cm

\noindent（2）使用边界与责任说明

\noindent 本学位论文/实践成果报告的核心观点、研究设计、实验方案、论证过程均由本人独立完成；人工智能生成内容均经过本人人工审核、修改与验证，并对其内容负责。本人对学位论文/实践成果报告全部内容的原创性、真实性、学术规范性承担全部责任，不存在由人工智能代写、伪造等学术不端行为。

\vskip 0.3cm

\noindent（3）承诺

\noindent 本人承诺人工智能工具在学位论文/实践成果报告中的使用符合国家相关法律法规、相关学术组织的规范及学校有关要求，不存在不当使用行为，愿意接受审查与监督，并独立承担相应责任和后果。

\vfill

\begin{flushright}
{作者签名}：$\underline{\qquad\qquad\qquad\qquad}$\quad

\vskip 0.5cm

{日\qquad 期}：\underline{\qquad\qquad} 年 \underline{\quad} 月 \underline{\quad} 日\quad
\end{flushright}
